\section{Conclusion}
\label{sec:conclusion}

We set out to combine two lines of work -- reconstruction attacks against released Rashomon sets of decision trees, and the effect of attacker background knowledge on reconstruction success -- into a single empirical study, and to push further into a gap neither line addresses: giving an individual reconstructed record a validated confidence value rather than treating the whole reconstruction as a single pass/fail event.

\paragraph{What we found.} Background knowledge of marital-status ($A$) and/or education ($B$) measurably reduces reconstruction error on the sensitive attribute \texttt{sex}, with a clean and consistent ordering $A{+}B < A \approx B < \text{None}$ that, at 350 seeds, survives Holm-Bonferroni-corrected per-$m$ significance testing for $A{+}B$ at all 5 released-tree counts, for $B$ at 3 of 5 ($m=10,25,50$), and for $A$ at 3 of 5 ($m=25,50,100$) (Section~\ref{sec:results}); an initial 200-seed pass had left $A$'s individual-$m$ significance unresolved, and the larger sample settled it in favor of a real, if smaller, effect. We looked for a structural explanation -- whether this effect is driven by released trees jointly splitting on $T$ and the known attribute -- and it mostly does not hold up under direct testing: a weak correlation is statistically detectable with enough data ($\rho\approx0.06$--$0.07$), but it explains well under 1\% of cell-to-cell variance, so it is not a practically useful predictor. We report that as a genuine (if nuanced) negative result rather than a nuisance to explain away. Finally, we built and validated a per-row confidence signal from solver agreement across repeated re-solves, showed it is predictive but not itself a calibrated probability, and then actually calibrated it with isotonic regression on a held-out seed split, cutting Brier score by 20--30\% and expected calibration error by a factor of 16--20$\times$ out of sample. Doubling the agreement scale's resolution (7 to 15 repeats) did not improve this; growing the held-out test set instead (6 to 10 seeds) did, most visibly at $m=100$ where calibrated ECE dropped from $0.072$ to $0.015$ -- confirming test-set size, not confidence-scale granularity, was the real bottleneck.

\paragraph{Why this is a meaningful contribution, not just a replication.} That releasing more trees from a Rashomon set lowers average reconstruction error is already established in prior work. What is new here is twofold: a direct, controlled ablation of \emph{which} background attributes an attacker knows (rather than just how many models are released), and -- to our knowledge, the first attempt in this literature -- a calibrated, held-out-validated confidence value attached to each individual reconstructed record, rather than a single aggregate error rate for the whole attack. An auditor or a data controller assessing real disclosure risk cares about both: which auxiliary information actually matters, and how much to trust any one specific reconstructed row, not just the average.

\paragraph{Scope and honesty about limitations.} This was explicitly built as an MVP within a three-week window, not an exhaustive study: a 13-column curated subset of a single dataset (UCI Adult), a single sensitive attribute, a single tree-selection strategy (enumeration order, no diversity-aware alternative), and the specific limitations already detailed in Section~\ref{sec:results-confidence}'s Limitations paragraph (pooled vs.\ per-$m$ test independence, solver timeouts at $m=100$, small agreement bins, held-out test-set size). None of the reported effects are overstated beyond what the corrected statistical tests support, and the one result that did not pan out (the joint-split mechanism) is reported as such rather than omitted.

\paragraph{Future work.} The most promising next steps are the ones that would sharpen what we already report rather than add new dimensions: a still-larger held-out calibration test set (10 seeds measurably helped over 6; more would likely help further, with diminishing but real returns) and a better mechanistic account of \emph{why} background knowledge helps beyond the joint-split proxy, which turned out to be statistically real but practically negligible (e.g.\ weighting leaves by how many training rows they actually constrain). Extending to additional sensitive attributes, datasets, or a differential-privacy defense would broaden the study but was judged out of scope for this project's timeline.
